\cotefile{14}
\scanpage{2}
\begin{flushright}
\fbox{\parbox{.50\linewidth}{\deleted{Les éléments de structure motiviques et leur interprétation en \unsure{termes} de groupes algébriques sur $\mathbb{Q}$}}}
\end{flushright}

\begin{center}
\uline{Structure d'effectivité}
\end{center}
% archive sheet number: 1
\scanpage{3}
\begin{center}Structure\end{center}

$4^{\circ})$ La donnée d'une sous-catégorie \inserted{\unsure{str.te|str.t|strict.} pleine} de $M$ (motifs effectifs), stable par $\otimes$, par sommes par \deleted{\unsure{sous et par}} sous-\deleted{\unsure{quo}}objets, telle que \inserted{où $\Lambda(-1)$ soit effectif et que} pour tt $E$, $\exists$ entier $n$ tel que $E(-n)=E\otimes\Lambda(-1)^{\otimes n}$ soit ``effectif''. Une telle structure sur $G_{\deleted{\unsure{C|\mathbb Q}}}$ en définit une autre sur $G_C$, et réciproquement, donc un tel élément de structure peut être considéré comme équivalent à la donnée d'un \deleted{partie de \unsure{$\chi^{+}$|$X^+$} $K(G)$ \unsure{base, à savoir}} \deleted{cône \unsure{$\chi$} $K^+(G)$} cône $K^+(G)$ \deleted{l'anneau \illegible{1} des \unsure{caractères|contenu} de $G_{\unsure{C|\mathbb C}}$ $(\chi^++\chi^+\subset\chi^+)$ \unsure{$\chi$} stable par produit, tel que pour tout} dans \inserted{l'anneau des \unsure{car.} \unsure{rat/$\mathbb{Q}$|virt./$\mathbb Q$}} $K(G)$, contenant $0$, $\deleted{\unsure{1|E}}\varepsilon$ stable par multiplication, et
% slanted
\inserted{engendré par \deleted{\unsure{les}} élém. de la base canonique \deleted{\illegible{1}} (classes de repr. irréd sur $\mathbb{Q}$) qu'il contient.}
À priori, ce semble là un élément de structure faisant intervenir la structure sur $Q$. Il y a lieu de poser d'autres axiomes. \inserted{Par exemple,} on voudra que pour tt \deleted{\unsure{M|H}} $\deleted{\unsure{V}}E\in\operatorname{Ob}M^+$ $E$ soit à graduations \uuline{positives}.
% slanted
\margin{De plus on voudra que la structure envisagée provienne d'une structure de \deleted{\unsure{$=$|même}} type donné sur $G^\circ$, \deleted{\unsure{d'où}} \unsure{par} \deleted{\illegible{1}} \unsure{supposant|supprimant} $G$ \unsure{connexe}.}
\margin{\uline{Axiome} Pour tt \deleted{\unsure{entier|degré} \unsure{$n$}} $i$ il n'y a qu'un nb fini de types de motifs effectifs de poids $i$ \ldots\ \textcircled{\scriptsize ???} En tous cas, on veut que les motifs effectifs de poids $0$ soient tels que $G^\circ$ opère en \unsure{eux} \unsure{trivialem.t}.}

Un lien $G_{m\,C}\xrightarrow{i}G_C$ \deleted{positif} \uline{effectif} si pour tout $E\in\operatorname{Ob}M^+$, la graduation correspondante de $E_C$ est à
% archive sheet number: 2
\scanpage{4}
\mspage{IV.2}
if and only if $\unsure{\mathbb{L}|\mathscr{L}}$ is equivalent to $\handwritten{\mathscr{A}}$.

\uline{Proof:} Certainly if $\mathscr{L}$ is in $\handwritten{\mathscr{C}}_1(\mathrm{X},\handwritten{\mathscr{A}})$ then $\mathscr{L}$ is equivalent to $\handwritten{\mathscr{A}}$. Suppose $\mathscr{L}$ is equivalent to $\handwritten{\mathscr{A}}$. Then there are faithful $\handwritten{\mathscr{A}}$-modules $\handwritten{\mathscr{M}}_1$ and $\handwritten{\mathscr{M}}_2$, locally projective of finite type, such that
\[
\mathscr{L}\otimes_{\handwritten{\mathscr{A}}}\underline{\mathrm{Hom}}_{\handwritten{\mathscr{A}}}(\handwritten{\mathscr{M}}_1,\handwritten{\mathscr{M}}_1)\approx \underline{\mathrm{Hom}}_{\handwritten{\mathscr{A}}}(\handwritten{\mathscr{M}}_2,\handwritten{\mathscr{M}}_2).
\]

Now $\underline{\mathrm{Hom}}_{\handwritten{\mathscr{A}}}(\handwritten{\mathscr{M}}_1,\handwritten{\mathscr{M}}_1)$ is imbedded in $\mathscr{L}\otimes_{\handwritten{\mathscr{A}}}\underline{\mathrm{Hom}}_{\handwritten{\mathscr{A}}}(\handwritten{\mathscr{M}}_1,\handwritten{\mathscr{M}}_1)$, so $\handwritten{\mathscr{M}}_2$ is a module over $\underline{\mathrm{Hom}}_{\handwritten{\mathscr{A}}}(\handwritten{\mathscr{M}}_1,\handwritten{\mathscr{M}}_1)$, which we will call $\handwritten{\mathcal{M}}$.

Now we have a map of presheaves $\alpha'\colon \underline{\mathrm{Hom}}_{\handwritten{\mathcal{M}}}(\handwritten{\mathscr{M}}_1,\handwritten{\mathscr{M}}_2)\otimes_{\handwritten{\mathscr{A}}}^{\handwritten{\prime}}\handwritten{\mathscr{M}}_1\to \handwritten{\mathscr{M}}_2$, defined by
\[
\alpha'_{\mathrm U}(\{\mathrm{g}_{\mathrm V}\}_{\mathrm V\subset\mathrm U}\}\otimes\mathrm{m}_1)=\mathrm{g}_{\mathrm U}(\mathrm{m}_1)
\]
for $\{\mathrm{g}_{\mathrm V}\}_{\mathrm V\subset\mathrm U}$ in $\mathrm{Hom}_{\handwritten{\mathcal{M}}|\mathrm U}(\handwritten{\mathscr{M}}_1|\mathrm U,\handwritten{\mathscr{M}}_2|\mathrm U)$ and $\mathrm{m}_1$ in $\handwritten{\mathscr{M}}_1(\mathrm U)$. Since $\handwritten{\mathscr{M}}_1$ has finite presentation over $\handwritten{\mathcal{M}}$ (Theorem II.7), we know from [5, Theorem A.6]that the induced map of sheaves, $\alpha$, is an isomorphism. Let $\underline{\mathrm{Hom}}_{\handwritten{\mathcal{M}}}(\handwritten{\mathscr{M}}_1,\handwritten{\mathscr{M}}_2)=\handwritten{\mathscr{M}}_3$. Since $\handwritten{\mathscr{M}}_1$ is locally projective of finite type as a $\handwritten{\mathcal{M}}$-module (by Theorem II.7) and $\handwritten{\mathscr{M}}_2$ is of finite type as a $\handwritten{\mathcal{M}}$-module, $\handwritten{\mathscr{M}}_3$ is of finite type as a $\handwritten{\mathcal{M}}$-module, and therefore is of finite type as an $\handwritten{\mathscr{A}}$-module. Thus, from Theorem II. 9, $\handwritten{\mathscr{M}}_3$ is locally projective of finite type
% archive sheet number: 3
\scanpage{5}
degrés positifs. On voit que si $i$ est effectif, il en est de \unsure{$\overline{\mathrm m}$|même} de ses conjugués par éléments de $G(C)$, et de ses conjugués par autom. de $C$, en particulier de $\bar i$. \deleted{\unsure{De}}D'autre part, pour $i$ donné, la sous-catégorie \inserted{pleine} $P(i)$ de $M$ \unsure{formé|formée} des $E$ tels que $E_C$ soit à $i$-graduation positive est \inserted{str. pleine,} stable par $\deleted{\otimes}$ sommes, par \deleted{\unsure{les}} sous-objet, par $\otimes$, contient \inserted{$0$, et} $\Lambda(-1)$ si on suppose que $\deleted{\varepsilon i(\lambda)=\lambda^\nu\ \text{avec }\nu}\deleted{\langle i,\varepsilon\rangle}\inserted{\varepsilon(i)}\geq 0$, et satisfait la condition que $\forall E\in\operatorname{Ob}M$, $\exists n$ tel que $E(\unsure{-n|n})\in\operatorname{Ob}\deleted{\unsure{M}}P(i)$, si on suppose $\boxed{\deleted{\langle i,\varepsilon\rangle}\inserted{\varepsilon(i)}>0}$ [N.B. on désigne par $\langle i,\varepsilon\rangle$ l'entier $\nu$ tel que $\varepsilon(i(\lambda))=\lambda^\nu$]. \deleted{Si l'on pose \unsure{égalem.t|engendré} pour les conjugués $\operatorname{int}(g)\circ i$ et les conjugués $\,{}^{\sigma}i$ $(g\in G(C),\ i\in$} \inserted{Prenant des intersections} des sous-catégories $P(i)$ on trouve des $M^+$ satisfaisant aux axiomes considérés, sauf peut être au dernier sur la graduation. \inserted{Par intersection de tels \unsure{jouets|fonct's|foncteurs}, itou \ldots}

Les $i$ pour lesquels, $M^+$ étant donné,
% archive sheet number: 4
\scanpage{6}
\begin{center}
IV. \uline{The Brauer group of a ringed space.}
\end{center}

The foregoing theorems enable us to define the Brauer group of a ringed space $(\mathrm{X},\handwritten{\mathscr{A}})$, where $\handwritten{\mathscr{A}}$ is a sheaf of commutative rings with unit.

We consider the category $\handwritten{\mathscr{C}}_1(\mathrm{X},\handwritten{\mathscr{A}})$ of locally separable $\handwritten{\mathscr{A}}$-algebras, with arrows the $\handwritten{\mathscr{A}}$-algebra morphisms. $\handwritten{\mathscr{C}}_1(\mathrm{X},\handwritten{\mathscr{A}})$ is closed under sheaf tensor product by Theorem III.5. We define $\handwritten{\mathscr{C}}'_1(\mathrm{X},\handwritten{\mathscr{A}})$ to be the full subcategory of $\handwritten{\mathscr{C}}_1(\mathrm{X},\handwritten{\mathscr{A}})$ consisting of those objects which are isomorphic to the algebras of the form $\underline{\mathrm{Hom}}_{\handwritten{\mathscr{A}}}(\handwritten{\mathscr{M}},\handwritten{\mathscr{M}})$, where $\handwritten{\mathscr{M}}$ is a faithful locally projective $\handwritten{\mathscr{A}}$-module of finite type (Theorem III,4). From Theorem II.4, we know that $\handwritten{\mathscr{C}}'_1(\mathrm{X},\handwritten{\mathscr{A}})$ is closed under tensor product.

We define an equivalence relation in $\handwritten{\mathscr{C}}_1(\mathrm{X},\handwritten{\mathscr{A}})$ as follows : $\mathscr{L}_1$ and $\mathscr{L}_2$ in $\handwritten{\mathscr{C}}_1(\mathrm{X},\handwritten{\mathscr{A}})$ are equivalent if and only if there exist $\handwritten{\mathcal{M}}_1$ and $\handwritten{\mathcal{M}}_2$ in $\handwritten{\mathscr{C}}'_1(\mathrm{X},\handwritten{\mathscr{A}})$ such that $\mathscr{L}_1\otimes_{\handwritten{\mathscr{A}}}\handwritten{\mathcal{M}}_1\approx \mathscr{L}_2\otimes_{\handwritten{\mathscr{A}}}\handwritten{\mathcal{M}}_2$.

The set of equivalence classes, $\handwritten{\mathfrak{B}}(\mathrm X,\handwritten{\mathscr{A}})$ is an abelian group under tensor product, with unit $\{\handwritten{\mathscr{A}}\}$ the class of $\handwritten{\mathscr{A}}$. If $\mathscr{L}$ is in $\handwritten{\mathscr{C}}_1(\mathrm{X},\handwritten{\mathscr{A}})$, then the inverse of $\{\mathscr{L}\}$ is $\{\mathscr{L}^{\unsure{\mathrm{o}|0}}\}$, since $\eta\colon \mathscr{L}\otimes_{\handwritten{\mathscr{A}}}\mathscr{L}^{\mathrm{o}}\to \underline{\mathrm{Hom}}_{\handwritten{\mathscr{A}}}(\mathscr{L},\mathscr{L})$ is an isomorphism.

THEOREM IV.1. If $\mathscr{L}$ is in $\handwritten{\mathscr{C}}_1(\mathrm{X},\handwritten{\mathscr{A}})$, then $\mathscr{L}$ is in $\handwritten{\mathscr{C}}'_1(\mathrm{X},\handwritten{\mathscr{A}})$
% archive sheet number: 5
\scanpage{7}
on a $M^+\subset P(i)$, \deleted{\unsure{peuvent|forment}} correspondent à certaines classes de conjugaison de liens. Si $T$ est un tore maximal de $G$, et $\Gamma=\operatorname{Hom}(G_{m\,C},T_C)$, $W=(\operatorname{Norm}(T)/T)(C)$, $W'=(\underline{\mathrm{Aut}}(G,T)/T)(C)$, de sorte que $W\subset W'$, $W'/W=\underline{\mathrm{Aut\ ext}}(G)(C)$, $W'$ opère sur $\Gamma$, et on obtient un \uline{cône} \deleted{de} $\Gamma^+$ de $\Gamma$, invariant par $W$ (mais peut-être pas par $W'$ !] contenant $j$, contenu dans l'hyperplan $\varepsilon\geq 0$ ; \deleted{ce cône étant} ce cône correspond donc à une partie ``conique'' de $\Gamma'=\Gamma/W$. Or $\Gamma'=\Gamma/W$
\[
\begin{aligned}
&\inserted{=[\underline{\mathrm{Hom}}(G_m,T)/(\operatorname{Norm}_G(T)/T)](C)\unsure{=|\Rightarrow}}\\
&=[\underline{\mathrm{Hom}}(G_m,G)/G](C)
\end{aligned}
\]
est formé des pts \deleted{\unsure{rat. sur}} \inserted{à val. ds} $C$ d'un schéma $\deleted{\unsure{E}}S$ non ramifié sur $Q$, et à ce titre le groupe fondamental $\pi_1(\mathbb{Q},C)$ opère sur \unsure{!} $\deleted{\unsure{S}}\Gamma=S(C)=S(\overline{\mathbb{Q}})$, d'ailleurs via un quotient fini de $\pi_1$. Ceci posé, le cône envisagé \deleted{contient} est stable par $\pi_1$.
% archive sheet number: 6
\scanpage{8}
\mspage{III.7}
of $\Gamma$, then $\Gamma^\Lambda=\mathrm{Hom}_{\Lambda^{\mathrm e}}(\Lambda,\Gamma)$ is called the \uline{commutant of $\Lambda$ in $\Gamma$} and is a subring of $\Gamma$. In [4, Theorem 3.5] the following is shown : if $\Gamma$ is central separable over $\mathrm A$, and $\Lambda$ and $\Omega$ are subalgebras of $\Gamma$ such that $\Lambda\otimes_{\handwritten{A}}\Omega\approx\Gamma$ under the mapping $\mathrm{x}\otimes\mathrm{y}\longrightarrow\mathrm{xy}$, then $\Lambda$ is the commutant of $\Omega$ and $\Omega$ is the commutant of $\Lambda$ in $\Gamma$.

Analogously, we can define the commutant of a subalgebra of sheaves. If $\handwritten{\mathscr{G}}$ is a sheaf of rings with center $\handwritten{\mathscr{A}}$, and $\mathscr{L}$ is an $\handwritten{\mathscr{A}}$-subalgebra of $\handwritten{\mathscr{G}}$, we define the commutant of $\mathscr{L}$ in $\handwritten{\mathscr{G}}$ to be the sheaf $\underline{\mathrm{Hom}}_{\mathscr{L}^{\mathrm e}}(\mathscr{L},\handwritten{\mathscr{G}})$, which is easily seen to be a subsheaf of $\handwritten{\mathscr{G}}$.

THEOREM III.7. Suppose $(\mathscr{L},\lambda)$ and $(\handwritten{\mathcal{M}},\unsure{\omega|\mathrm w})$ are two sheaves of algebras, locally separable over $(\handwritten{\mathscr{A}},\alpha)$. Then $\mathscr{L}$ is isomorphic to the commutant of $\handwritten{\mathcal{M}}$ in $\mathscr{L}\otimes_{\handwritten{\mathscr{A}}}\handwritten{\mathcal{M}}$, and $\handwritten{\mathcal{M}}$ is isomorphic to the commutant of $\mathscr{L}$.

\uline{Proof:} We will define an $\handwritten{\mathscr{A}}$-map $\mathrm{g}\colon \mathscr{L}\to \underline{\mathrm{Hom}}_{\handwritten{\mathcal{M}}^{\mathrm e}}(\handwritten{\mathcal{M}},\mathscr{L}\otimes_{\handwritten{\mathscr{A}}}\handwritten{\mathcal{M}})$.

Let $\mathrm U$ be an open set in $\mathrm X$, and $\mathrm p$ an element of $\mathscr{L}(\mathrm U)$. Define $\mathrm{g}'_{\mathrm U}(\mathrm p)$ to be the map of presheaves in $\mathrm{Hom}_{\handwritten{\mathcal{M}}^{\mathrm e\handwritten{'}}|\mathrm U}(\handwritten{\mathcal{M}}|\mathrm U,\mathscr{L}|\mathrm U\otimes_{\handwritten{\mathscr{A}}|\mathrm U}^{\handwritten{\prime}}\handwritten{\mathcal{M}}|\mathrm U)$ such that $\mathrm{g}'_{\mathrm U}(\mathrm p)_{\mathrm V}(\mathrm q)=\lambda_{\mathrm V}^{\mathrm U}(\mathrm p)\otimes\mathrm q$ for any open set $\mathrm V$ contained in $\mathrm U$ and $\mathrm q$ in $\handwritten{\mathcal{M}}(\mathrm V)$. Let $\overset{\handwritten{\rule{2.4em}{0.4pt}}}{\mathrm{g}'_{\mathrm U}(\mathrm p)}$ be the sheaf map induced by $\mathrm{g}'_{\mathrm U}(\mathrm p)$. Then $\mathrm{g}_{\mathrm U}(\mathrm p)$ will be defined as $\overset{\handwritten{\rule{2.4em}{0.4pt}}}{\mathrm{g}'_{\mathrm U}(\mathrm p)}$ for every open set $\mathrm U$ in $\mathrm X$ and $\mathrm p$ in $\mathscr{L}(\mathrm U)$.

Since $\handwritten{\mathcal{M}}$ is an $\handwritten{\mathscr{A}}$-module of finite type, it is finitely presented
% archive sheet number: 7
\scanpage{9}
\deleted{À tout ``cône'' \unsure{stable|stabl}}

Si nous nous intéressons aux structures d'effectivité \inserted{qui peuvent se} définir \deleted{par} en termes de \deleted{\unsure{certain}} une famille de $\underline{i}$ \deleted{\unsure{pour}}comme intersection des $P(i)$, \deleted{\unsure{on a}} la donnée d'une telle structure équivaut à celle d'un ``cône'' de $\Gamma'$, invariant par $\pi_1$, contenu \inserted{``\uline{strictement}''} dans celui défini par $\varepsilon$,
% slanted
\margin{\unsure{``''|!!!} i.e. $\varepsilon$ est \unsure{``}\uline{intérieur} au cône polaire !}
contenant la classe de $i$, satisfaisant à une condition de \deleted{\unsure{non} aplatissement} saturation que nous ne précisons pas ici (est-elle automatiquement satisfaite ?).

Le cas le plus intéressant est celui \deleted{\unsure{de}} \inserted{où} $M^+$ \deleted{\unsure{$Q$|$\otimes$}} est défini comme $P(i_1)$, où $i_1$ est comme au $3^{\circ}$. \deleted{\unsure{En particulier},} \inserted{[N.B.} le cône \uline{invariant} engendré contient bien $j$, et est contenu dans le $\frac12$ espace défini par $\varepsilon$, puisque $i_1\bar i_1=j$, $\varepsilon i_1=\mathrm{id}$]. On trouve ainsi une application
\[
I(G)\longrightarrow\deleted{\unsure{\mathrm{Pol}}}\mathrm{Eff}(G)
\]
% archive sheet number: 8
\scanpage{10}
\mspage{III.6}
THEOREM III.6. If $\mathscr{L}$ is a locally separable $\handwritten{\mathscr{A}}$-algebra, and $\handwritten{\mathscr{R}}$ is a commutative $\handwritten{\mathscr{A}}$-algebra, then $\mathscr{L}\otimes_{\handwritten{\mathscr{A}}}\handwritten{\mathscr{R}}$ is a locally separable $\handwritten{\mathscr{R}}$-algebra.

\uline{Proof:} We know\handwritten{\unsure{,|$\swarrow$}} from [4, Corollary 1.6] that $(\mathscr{L}\otimes_{\handwritten{\mathscr{A}}}\handwritten{\mathscr{R}})_{\mathrm x}=\mathscr{L}_{\mathrm x}\otimes_{\handwritten{\mathscr{A}}_{\mathrm x}}\handwritten{\mathscr{R}}_{\mathrm x}$ is a central separable $\handwritten{\mathscr{R}}_{\mathrm x}$-algebra for every $\mathrm x$ in $\mathrm X$. Thus we need only show that $\mathscr{L}\otimes_{\handwritten{\mathscr{A}}}\handwritten{\mathscr{R}}$ has finite presentation as an $\handwritten{\mathscr{R}}$-algebra. For every $\mathrm x$ in $\mathrm X$ there is a neighborhood $\mathrm U$ of $\mathrm X$ and an exact sequence
\diagram{\begin{tikzcd}[ampersand replacement=\&,column sep=small]
0\handwritten{\arrow[r]}\& \handwritten{\mathscr{P}}\handwritten{\arrow[r]}\& \overset{\mathrm n}{\oplus}\,\handwritten{\mathscr{A}}|\mathrm U\handwritten{\arrow[r]}\& \mathscr{L}|\mathrm U\handwritten{\arrow[r]}\& 0
\end{tikzcd}}
where $\handwritten{\mathscr{P}}$ is an $\handwritten{\mathscr{A}}|\mathrm U$-module of finite type. Then the sequence $\handwritten{\mathscr{P}}\otimes_{\handwritten{\mathscr{A}}|\mathrm U}\handwritten{\mathscr{R}}|\mathrm U\handwritten{\longrightarrow}\overset{\mathrm n}{\oplus}\,\handwritten{\mathscr{A}}|\mathrm U\otimes_{\handwritten{\mathscr{A}}|\mathrm U}\handwritten{\mathscr{R}}|\mathrm U\handwritten{\longrightarrow} \mathscr{L}|\mathrm U\otimes_{\handwritten{\mathscr{A}}|\mathrm U}\handwritten{\mathscr{R}}|\mathrm U\handwritten{\longrightarrow} 0$ is exact, and $\handwritten{\mathscr{P}}\otimes_{\handwritten{\mathscr{A}}|\mathrm U}\handwritten{\mathscr{R}}|\mathrm U$ is an $\handwritten{\mathscr{R}}|\mathrm U$-module of finite type. Thus $\mathscr{L}\otimes_{\handwritten{\mathscr{A}}}\handwritten{\mathscr{R}}$ is finitely presented as an $\handwritten{\mathscr{R}}$-module, and therefore locally separable.

We recall now the notion of the commutant of a subring in a given ring, as described in [4]. This notion, generalized to sheaves, will be of use in the next section.

For any ring $\Lambda$ with center $\mathrm A$, and any two-sided $\Lambda$-module $\mathrm M$, $\mathrm{Hom}_{\Lambda^{\mathrm e}}(\Lambda,\mathrm M)$ is naturally isomorphic as an $\mathrm A$-module with $\mathrm M^\Lambda$, the $\mathrm A$-submodule of $\mathrm M$ consisting of all $\mathrm m$ in $\mathrm M$ such that $\mathrm{xm}=\mathrm{mx}$ for all $\mathrm x$ in $\Lambda$. If $\Gamma$ is a ring having center $\mathrm A$, and $\Lambda$ is an $\mathrm A$-subalgebra
% archive sheet number: 9
\scanpage{11}
\deleted{notant que}

en particulier, tenant compte de $3^\circ$,
\[
I^{\unsure{\alpha|*|+}}(G)\longrightarrow \mathrm{Pol}(G)\times\mathrm{Eff}(G)
\]
\margin{\unsure{''} les $i:G_{m\,C}\to G_C$ définissant une polarisation, modulo \unsure{action} de $G(\mathbb R)$ \ldots}

\uline{Questions} Cette dernière application, ou $I(G)\unsure{-|\to}\mathrm{Eff}(G)$, injectives ?

Les structures $\in\mathrm{Eff}(G)$ provenant de la théorie des motifs peuvent elles tjrs se décrire par un $i_1$ ? i.e. sont elles dans $\operatorname{Im}I(G)$ ?
\margin{N.B. Le fait qu'une structure d'effectivité puisse se décrire par un élément de $I(G)$ resp $I^*(G)$ est une propriété intrinsèque à la catégorie de motifs $M$, indép. du foncteur qui la réalise. ! !}

(Regarder ce qu'on obtient en tordant par un torseur). Examiner aussi, pour ces questions et celles posées dans $3^\circ$, le cas où $G$ avec ses 4 structures \unsure{provient|provenant} d'un motif sur schéma \inserted{connexe} $S$ muni d'un $\xi\in S(\mathbb{Q})$, en utilisant le foncteur de De-Rham \ldots\ ? ]

Regarder $G$ commutatif ! $G=\unsure{\mathrm{Sl}|\mathrm{SL}|\mathrm{Gal}}(\unsure{2|n|?|\overline{\mathbb Q}},\mathbb{Q})$ !

Des propriétés arithmétiques plus profondes pourront être sans doute formulées par l'introduction des éléments de \unsure{frobenius|Frobenius|filtrations} et l'interprétation des structures $3^\circ$ et $4^\circ$ en termes de ces derniers.
% archive sheet number: 10
\scanpage{12}
\mspage{III.5}
the fact that $\mathrm M\otimes_{\mathrm A}\mathrm{Hom}_{\mathrm A}(\mathrm M,\mathrm A)$ is projective as a left $\mathrm W^{\mathrm e}$-module. Since $\mu$ is a left $\mathrm W^{\mathrm e}$-isomorphism, we see that $\mathrm W$ is projective as a left $\mathrm W^{\mathrm e}$ module, i. e., $\mathrm W$ is separable over $\mathrm A$.

THEOREM III.4. If $\handwritten{\mathscr{M}}$ is a faithful locally projective sheaf of modules of finite type over $\handwritten{\mathscr{A}}$, then $\underline{\mathrm{Hom}}_{\handwritten{\mathscr{A}}}(\handwritten{\mathscr{M}},\handwritten{\mathscr{M}})$ is a locally separable $\handwritten{\mathscr{A}}$-algebra.

\uline{Proof:} We know from Theorem II.2. that $\underline{\mathrm{Hom}}_{\handwritten{\mathscr{A}}}(\handwritten{\mathscr{M}},\handwritten{\mathscr{M}})=\Omega$ is a locally projective module of finite type. Since $\handwritten{\mathscr{M}}$ has finite presentation as an $\handwritten{\mathscr{A}}$-module, $\Omega_{\mathrm x}=\mathrm{Hom}_{\handwritten{\mathscr{A}}_{\mathrm x}}(\handwritten{\mathscr{M}}_{\mathrm x},\handwritten{\mathscr{M}}_{\mathrm x})$ for every $\mathrm x$ in $\mathrm X$. Since $\handwritten{\mathscr{M}}_{\mathrm x}$ is a finitely generated faithful $\handwritten{\mathscr{A}}_{\mathrm x}$-module for every $\mathrm x$, $\Omega_{\mathrm x}$ is a central separable $\handwritten{\mathscr{A}}_{\mathrm x}$-algebra for every $\mathrm x$. Thus by Theorem III.2, $\Omega$ is a locally separable $\handwritten{\mathscr{A}}$-algebra.

THEOREM III.5. If $\mathscr{L}_1$ and $\mathscr{L}_2$ are locally separable over $\handwritten{\mathscr{A}}$, then $\mathscr{L}_1\otimes_{\handwritten{\mathscr{A}}}\mathscr{L}_2$ is locally separable.

\uline{Proof:} We know that for every $\mathrm x$ in $\mathrm X$, $(\mathscr{L}_1\otimes_{\handwritten{\mathscr{A}}}\mathscr{L}_2)_{\mathrm x}\approx \mathscr{L}_{1_{\mathrm x}}\otimes_{\handwritten{\mathscr{A}}_{\mathrm x}}\mathscr{L}_{2_{\mathrm x}}$. Since $\mathscr{L}_{1_{\mathrm x}}$ and $\mathscr{L}_{2_{\mathrm x}}$ are central separable over $\handwritten{\mathscr{A}}_{\mathrm x}$, we have, from [4, Proposition 1.5] the fact that $\mathscr{L}_{1_{\mathrm x}}\otimes_{\handwritten{\mathscr{A}}_{\mathrm x}}\mathscr{L}_{2_{\mathrm x}}$ is central separable over $\handwritten{\mathscr{A}}_{\mathrm x}$ for every $\mathrm x$ in $\mathrm X$. Thus Theorem III.2 tells us that $\mathscr{L}_1\otimes_{\handwritten{\mathscr{A}}}\mathscr{L}_2$ is a locally separable $\handwritten{\mathscr{A}}$-algebra.
% archive sheet number: 11
\scanpage{13}
\deleted{Considérons un point $\xi\in S(k)$, \inserted{$k$ corps de car. $0$,} il définit un foncteur fibre $H^{**}_\xi$ (Hodge) et $H^*_\xi$ (De Rham) à valeurs dans la catégorie des vectoriels sur $k$.}

\begin{center}\rule{3cm}{0.4pt}\end{center}

Cas d'un \uline{tore} $T$ sur $\mathbb Q$.

Défini par un $\pi$-module admissible, libre de t.f. sur $\mathbb Z$, soit $M$, $(\pi=\pi_1(\mathbb Q,C))$

On aura $\varepsilon\in M^\pi$, $i\deleted{\unsure{\in}}\in\check M^\pi$, $\langle i,\varepsilon\rangle=2$.

Les modules sur $T$ correspondent aux éléments invariants par $\pi$ de $\mathbb Z(M)^+$ (comb. lin. à coeff. $\geq 0$)

\deleted{Donnons nous \unsure{le|$\varepsilon$} un}

La donnée d'une structure \deleted{\unsure{positive}} \inserted{de polarisation} ne dépend que du tore réel $T_{\mathbb R}$. Ce dernier est de la forme \deleted{\unsure{$T_2$|$T_{\mathbb R}\simeq$} $U\times S$,} \inserted{ext. d'un $S$ \unsure{splitté|split} par} $U$ compact, \inserted{i.e. $U(\mathbb R)$ cpct,} \deleted{$S$ \unsure{splitté|split},} et $\varepsilon$ est néc. triviale sur $U$. Donc on en déduit des structures $\varepsilon_S$ et $j_S$ sur $S$, et évidemment $S$ \uline{hérite} de structure de positivité. J'affirme qu'on doit avoir $\dim S=1$, en d'autres termes, que \uline{$\operatorname{Ker}\varepsilon$ est compact}.

En effet, \deleted{si on avait $\dim S\geq 2$, on \deleted{pourrait} peut \deleted{trouver un} se ramener (divisant par \deleted{\unsure{qq}} un sous-tore de codim 1 de $\operatorname{Ker}\varepsilon$) que la dim est deux. \deleted{Soit} On \deleted{\unsure{aura}} peut écrire $S=G_{m\,\mathbb R}\times G_{m\,\mathbb R}$, avec $i(\deleted{\unsure{t|1}}\lambda)=(\unsure{1|0|\lambda^\beta},\lambda^\alpha)$, $\varepsilon(\lambda,\mu)=\mu^{\unsure{\beta}}$, où $\alpha,\beta$ entiers $\geq 1$, $\alpha\beta=2$ (donc $\alpha\unsure{=|,\ }1$, $\beta=2$, ou $\alpha=2$, $\beta=1$).}
% archive sheet number: 12
\scanpage{14}
\mspage{III.4}
Using this fact and Theorem III.1, we see immediately that the following theorem is true :

THEOREMIII.2. $\mathscr{L}$ is locally separable over $\handwritten{\mathscr{A}}$ if and only if $\mathscr{L}$ is finitely presented as an $\handwritten{\mathscr{A}}$-module and $\mathscr{L}_{\mathrm x}$ is a central separable $\handwritten{\mathscr{A}}_{\mathrm x}$-algebra for every $\mathrm x$ in $\mathrm X$.

It is clear that if $\mathscr{L}$ is locally separable over $\handwritten{\mathscr{A}}$ and $\handwritten{\mathscr{M}}$ is a module over $\mathscr{L}$, locally projective of finite type, then $\handwritten{\mathscr{M}}$ is locally projective of finite type over $\handwritten{\mathscr{A}}$. The following lemma concerns ordinary rings and modules (not sheaves).

LEMMA III.3. If $\mathrm A$ is a commutative ring with unit and $\mathrm M$ is a finitely generated projective $\mathrm A$-module, then $\mathrm{Hom}_{\mathrm A}(\mathrm M,\mathrm M)$ is separable over $\mathrm A$.

\uline{Proof:} Let $\mathrm{Hom}_{\mathrm A}(\mathrm M,\mathrm M)=\mathrm W$. From [5, Proposition A.1] we see that since $\mathrm M$ is a finitely generated projective $\mathrm A$-module, the \inserted{mapping $\mu\colon\mathrm{Hom}_{\mathrm A}(\mathrm M,\mathrm A)\otimes_{\mathrm A}\mathrm M\to\mathrm W$, defined by} $\mu(\mathrm f\otimes\mathrm m)(\mathrm n)=\mathrm f(\mathrm n)(\mathrm m)$ is a two-sided $\mathrm W$-isomorphism and therefore a left $\mathrm W^{\mathrm e}$-isomorphism. From [5, Theorem A.2] we know that $\mathrm{Hom}_{\mathrm A}(\mathrm M,\mathrm A)$ is finitely generated and projective as a right $\mathrm W$-module, and therefore as a left $\mathrm W^{\unsure{\mathrm{o}|0}}$-module, and from [\unsure{5|3}, Proposition A.3] we know that $\mathrm M$ is finitely generated and projective as a left $\mathrm W$-module. Thus from [\handwritten{7}, p.165, Proposition 2.3] we have
% archive sheet number: 13
\scanpage{15}
Détermination du groupe des caractères rationnels sur $\mathbb Q$

$T_{\mathbb C}=D_{\mathbb C}(M)$, $M$ réseau où $\pi_1$ opère \inserted{via un \unsure{gpe|groupe} quotient fini $\pi$ opérant fidèlement} $(\pi_1=\pi_1(\mathbb Q,\mathbb C))$

$\mathbb Z(M)$ anneau des caractères \inserted{virtuels} de $T_{\mathbb C}$, $\pi$ y opère.

$K(T)=\mathbb Z(\deleted{\unsure{N|M}}M)^\pi$\qquad Anneau de type fini sur $\mathbb Z$.

[N.B. $\mathbb Z(M)$ est l'anneau affine de $S=\deleted{\unsure{\mathbb Z}}D_{\mathbb Z}(M)$, sur lequel $\pi$ opère par automorphismes de groupes. $K(T)$ est l'anneau affine de $D_{\mathbb Z}(M)/\pi$]. Les caractères \uline{effectifs} sont ceux qui prennent des valeurs $\geq 0$ sur les points de $S(\mathbb R)^0$ ?? \margin{$\lambda_1^{n_1}\lambda_2^{n_2}\unsure{-|\cdots}\lambda_r^{n_r}$} Comment définir une structure d'effectivité \uline{motivique} ? Le procédé \unsure{standart|standard} est de \deleted{se} fixer un $i_1:G_{m\mathbb C}-T_{\mathbb C}$ i.e. un $i_1^*:M\unsure{-|\to}\mathbb Z$, avec les propriétés habituelles, \deleted{et} donnant donc $\mathbb Z(M)-\mathbb Z(\mathbb Z)$ d'où $G_{m\mathbb Z}\xrightarrow{\alpha}S$ (par \unsure{comp.|conj.} avec $\pi$....)

La condition est, \deleted{que} pour un $\chi\in K(T)\simeq\mathbb Z(S)^\pi$, que $\chi\circ\alpha\simeq\deleted{\unsure{\mathbb Z|G}}$ anneau affine de $G_{m\mathbb Z}$ soit une fonction \uline{polynômiale} i.e. \deleted{restant} finie sur $G_{m\mathbb Z}(\mathbb R)=\mathbb R^*$ pour $t\to0$.

Considérons les caractères de motifs effectifs de poids donné \unsure{$s$|$p$|$\rho$}, prouvons qu'il n'y en a qu'un nb fini. Cela \inserted{\uuline{ne}} \uuline{\unsure{tient|peut|vient}} \inserted{\uuline{pas}} simplement \unsure{du|en} \unsure{calcul|critère} \unsure{arithmétique|quantième|quand on passe de} $\mathbb Q$ par $\mathbb R$, \deleted{Alors} car si $T$ est de dim $\geq3$, la relation de finitude correspondante est \uline{fausse}.

% archive sheet number: 14
